% Generated draft: requires mathematical and source review.
\documentclass[11pt]{article}
\usepackage[margin=25mm]{geometry}
\usepackage{fontspec}
\setmainfont{Latin Modern Roman}
\usepackage{amsmath,amssymb,mathtools,unicode-math}
\setmathfont{Latin Modern Math}
\usepackage{enumitem,needspace,xurl,hyperref}
\setlength{\parindent}{0pt}
\setlength{\parskip}{5pt}
\setlength{\emergencystretch}{3em}
\newcommand{\N}{\mathbb{N}}
\newcommand{\Z}{\mathbb{Z}}
\newcommand{\Q}{\mathbb{Q}}
\newcommand{\R}{\mathbb{R}}
\newcommand{\C}{\mathbb{C}}
\newcommand{\F}{\mathbb{F}}
\newcommand{\A}{\mathbb{A}}
\newcommand{\PP}{\mathbb{P}}
\newcommand{\E}{\mathbb{E}}
\newcommand{\eps}{\varepsilon}
\newcommand{\dd}{\,\mathrm{d}}
\begin{document}
\section*{UnsolvedMath ID: 11300004}
\textbf{Source problem code:} catalogue entry
\subsection{Definitions and mathematical statement}
A knot \(K\subset\R^3\) is the image of a topological embedding \(S^1\hookrightarrow\R^3\). It is tame if it is ambient isotopic to a polygonal knot, and wild otherwise. Define its quadrisecant set by
\[Q(K)=\{\ell\subset\R^3:\ell\text{ is an affine line and }|\ell\cap K|\geq4\}.\]
Lines are counted without orientation, and four distinct intersection points suffice; a line meeting the knot in more points is also included and counts as one line.

The question attributed to Greg Kuperberg is
\[\forall K\text{ wild},\qquad |Q(K)|=\infty.\]
Equivalently, for every such \(K\) and every integer \(N\geq1\), there exist pairwise distinct lines \(\ell_1,\ldots,\ell_N\), each meeting \(K\) in at least four distinct points. Infinitely many point quadruples lying on one line would not suffice. No smoothness, generic-position, transversality, cyclic-order or alternating-quadrisecant condition is imposed.

The research summary discusses classical nontrivial tame-knot existence results attributed to Pannwitz and possible finite quadrisecant counts, but those are context for a different hypothesis and conclusion. Its wording about a generic isotopy-invariant ``Q-number'' has no definition or supporting reference here and is not adopted as a property of the raw set \(Q(K)\). The summary's tentative nontrivial/knotted wording is not added as an extra restriction on the actual statement, which quantifies over every wild knot.
\subsection{Short English statement}
Must every wild embedded circle in three-space have infinitely many distinct affine lines each meeting it in at least four distinct points? The question does not require exactly four points or generic intersections.
\subsection{Sources}
\begin{itemize}
\item Local UnsolvedMath registry, permanent ID \(11300004\), exact source code \texttt{catalogue entry}; proposed by Greg Kuperberg in the supplied record.
\item Wikipedia, \emph{List of unsolved problems in mathematics}, Topology, ``Quadrisecants of wild knots,'' revision 1366636767: \url{https://en.wikipedia.org/w/index.php?title=List_of_unsolved_problems_in_mathematics&oldid=1366636767#Topology}. Registry access 2026-07-29; MediaWiki-source extraction.
\item The summary mentions Pannwitz and Kuperberg's 1994/1996 \emph{Quadrisecants of knots} context but gives no full primary bibliographic entry. No authoritative primary citation is supplied in this record.
\end{itemize}
\end{document}
